% Corsin Week 3 -- full exercise sheet (complete; this week has no problems tagged
% important, so content/week-03.tex carries none of them inline). NOT \input by main.tex --
% reference source only.
\exercisesheet{3}

\textit{Statements quoted verbatim from \texttt{exercises/Ex3\_Analysis2\_eng.pdf} (assigned 2
March 2026, due 9 March 2026). Attribution: Prof. Joaquim Serra, D-MATH, ETH Z\"urich.}

\begin{exercise}[3.1 --- Lipschitz or not]
\label{ex:3.1}
Give an example of:
\begin{enumerate}[label=\textbf{(\alph*)}]
  \item A $\tfrac{1}{2}$-Lipschitz function $f : [0,1) \to [0,1)$ that has no fixed point.
  \item A map $f : \mathbb{R}^2 \to \mathbb{R}^2$ that has no fixed point, but is an isometry,
    i.e.\ $\|f(x) - f(x')\| = \|x - x'\|$ for all $x, x' \in \mathbb{R}^2$.
\end{enumerate}
\end{exercise}

\begin{exercise}[3.2 --- Problems at the origin]
\label{ex:3.2}
Show that there is no continuous function $f : \mathbb{R}^2 \to \mathbb{R}$ such that
$f(x,y) = xy/(x^2+y^2)$ for all $(x,y) \neq (0,0)$. On the other hand, show that there is exactly
one continuous function $g : \mathbb{R}^2 \to \mathbb{R}$ such that
$g(x,y) = xy/\sqrt{x^2+y^2}$ for all $(x,y) \neq (0,0)$.

\textit{Official hint:} if a function is continuous at $0$ and $x_k \to 0$ and $y_k \to 0$, then
$f(x_k)$ and $f(y_k)$ must converge to the same value.
\end{exercise}

\begin{exercise}[3.3 --- Open, closed, complete and compact]
\label{ex:3.3}
For each of the following subsets of $\mathbb{R}^N$ say whether they are open / closed / none /
compact (with respect to the standard Euclidean structure). Try to prove your assertions
``efficiently''.
\begin{enumerate}[label=\textbf{\arabic*.}]
  \item $E_1 := \{x \in \mathbb{R}^3 : 0 < x_2 \leq 2\}$
  \item $E_2 := \{x \in \mathbb{R}^n : \sin(\|x\|) \geq \tfrac{1}{4}\}$
  \item $E_3 := \bigcup_{n \geq 1}\{x \in \mathbb{R}^3 : x_1^4 - \tfrac{1}{n}x_3^2 - x_2^2 > \tfrac{1}{n},\ x_1 + x_2 < 6\}$
  \item $E_4 := \{(x,y) \in (\mathbb{R}^n \times \mathbb{R}^n) : x \cdot y > \tfrac{1}{2}\|x\|\|y\|\}$
  \item $E_5 := \{X \in \mathbb{R}^{n\times n} : \text{the matrix } X \text{ is invertible}\}$
  \item $E_6 := \{X \in \mathbb{R}^{n\times n} : \text{the matrix } X \text{ is symmetric}\}$
  \item $E_7 := \{X \in \mathbb{R}^{n\times n} : \text{the entries of } X \text{ are either } 0 \text{ or } 1\}$
  \item $E_8 := \{x \in \mathbb{R}^n : |x_i| \leq 6 \text{ for all } 1 \leq i \leq n\}$
  \item $E_9 := E_8 \times E_3 \subset \mathbb{R}^{n+3}$
  \item $E_{10} := E_2 \cap E_8 \subset \mathbb{R}^n$
\end{enumerate}
\end{exercise}

\begin{exercise}[3.4 --- Multiple choice (open, complete, compact)]
\label{ex:3.4}
Select all the statements below that are true.
\begin{enumerate}[label=\textbf{(\alph*)}]
  \item A nonempty open strict subset of $\mathbb{R}^n$ cannot be compact.
  \item A nonempty open strict subset of $\mathbb{R}^n$ cannot be complete.
  \item A complete subset of $\mathbb{R}^n$ contains all its accumulation points.
  \item Countable union of complete subsets of $\mathbb{R}^n$ is complete.
  \item A subset is closed in $\mathbb{R}^n$ if and only if it is complete as a metric space
    itself (with the distance inherited from $\mathbb{R}^n$).
\end{enumerate}
\end{exercise}

\begin{exercise}[3.5 --- Multiple choice (fixed points)]
\label{ex:3.5}
Select all the statements below that are true.
\begin{enumerate}[label=\textbf{(\alph*)}]
  \item If you spread a map of Zurich on your desk, then one point of the desk will coincide with
    its representation on the map (in an ideal world).
  \item If $f : [0,1] \to [0,2]$ is $\tfrac{1}{2}$-Lipschitz, then $f$ has a fixed point.
  \item If $f : [0,2] \to [0,1]$ is $\tfrac{1}{2}$-Lipschitz, then $f$ has a fixed point.
  \item If $f : [0,1] \cup [2,3] \to [0,1] \cup [2,3]$ is continuously differentiable with
    $|f'(x)| < 1$ for all $x \in [0,1]\cup[2,3]$, then it has a fixed point.
  \item \textbf{(*)} If $f : [0,1] \to [0,1]$ is differentiable with $|f'(x)| < 1$ for all
    $x \in (0,1)$, then it has a fixed point.
\end{enumerate}
\end{exercise}

\begin{exercise}[3.6 --- Multiple choice (uniform continuity)]
\label{ex:3.6}
Select all the statements below that are true.
\begin{enumerate}[label=\textbf{(\alph*)}]
  \item The function $(x,y) \mapsto x+y$ is uniformly continuous in $\mathbb{R}^2$.
  \item The function $(x,y) \mapsto xy$ is uniformly continuous in $\mathbb{R}^2$.
  \item The function $(x,y) \mapsto x+y$ is uniformly continuous in $[0,1]^2$.
  \item The function $(x,y) \mapsto xy$ is uniformly continuous in $[0,1]^2$.
\end{enumerate}
\end{exercise}

\begin{remark}
Corsin's third reason to care about compactness (\cref{sec:compactness_continued}) is precisely
that a continuous function on a compact space is uniformly continuous --- which settles
\cref{ex:3.6}(c) and (d) at once.
\end{remark}

\begin{exercise}[3.7 --- Cauchy--Schwarz]
\label{ex:3.7}
Let $V$ be a vector space over $\mathbb{R}$, let $\langle\cdot,\cdot\rangle$ be an inner product
on $V$, and let $\|\cdot\| : V \to \mathbb{R}$ be given by $\|v\| = \sqrt{\langle v,v\rangle}$.
Then the inequality
\begin{equation*}
|\langle v,w\rangle| \leq \|v\|\|w\| \tag{1}
\end{equation*}
holds for all $v, w \in V$. Furthermore, equality in (1) holds if and only if $v$ and $w$ are
linearly dependent.
\end{exercise}

\begin{remark}
Corsin devotes the whole second half of Friday (\cref{sec:cauchy_schwarz}) to this problem,
including the geometric reading and an explicit proof hint.
\end{remark}

\begin{exercise}[3.8 --- Inner product and norm]
\label{ex:3.8}
Let $V$ be a vector space over $\mathbb{R}$, let $\langle -,-\rangle$ be an inner product on $V$.
The map defined by
$$\|\cdot\| : V \to \mathbb{R}, \qquad \|v\| = \sqrt{\langle v,v\rangle}$$
satisfies the triangular inequality and is a norm.
\end{exercise}

\begin{exercise}[3.9 --- All norms are equivalent in $\mathbb{R}^n$]
\label{ex:3.9}
Let $|\cdot|$ denote the standard Euclidean norm in $\mathbb{R}^n$ and let
$f : \mathbb{R}^n \to [0,\infty)$ be another norm (that is, a function satisfying the properties
of \textit{Definition 9.89}).
\begin{enumerate}[label=\textbf{\arabic*.}]
  \item Expressing $x$ in a basis and using the ``abstract'' properties that $f$ must have, show
    that there is a constant $C_1 > 0$ such that $f(x) \leq C_1|x|$ for all $x \in \mathbb{R}^n$.
  \item Show that $f$ is continuous in $\mathbb{R}^n$ (with respect to the standard distance of
    $\mathbb{R}^n$!).
  \item Show that there is a number $c_2 > 0$ such that $f(x) \geq c_2$ for all $|x| = 1$.
  \item Conclude that $f(x) \geq c_2|x|$ for all $x \in \mathbb{R}^n$.
  \item Show that if $\tilde f$ is yet another norm, then there is $C > 0$ such that
    $C^{-1}f(x) \leq \tilde f(x) \leq Cf(x)$ for all $x \in \mathbb{R}^n$.
\end{enumerate}
\textit{Official hints:} \textbf{3.9.2} --- recall that from the definition it follows that
$|f(x) - f(y)| \leq f(x-y)$, and that Lipschitz functions are always continuous. \textbf{3.9.3}
--- apply the Weierstrass theorem to $f$ on $S := \{x \in \mathbb{R}^n : |x| = 1\}$ and use that
$f$, being a norm, is nondegenerate. \textbf{3.9.5} --- if $f$ is equivalent to $|\cdot|$ and
$\tilde f$ is equivalent to $|\cdot|$, it follows by transitivity that $f$ is equivalent to
$\tilde f$.
\end{exercise}

\begin{remark}
Hint 3.9.3 is exactly Corsin's reason \#2 for caring about compactness
(\cref{sec:compactness_continued}) --- the Weierstrass/extreme value theorem on the compact unit
sphere.
\end{remark}

\begin{exercise}[3.10 --- Hilbert--Schmidt norm of the composition]
\label{ex:3.10}
Take two linear functions $\varphi : \mathbb{R}^d \to \mathbb{R}^n$ and
$\psi : \mathbb{R}^n \to \mathbb{R}^m$, and denote with $\Phi$, $\Psi$ the matrices that represent
them in the canonical bases. Recall that the linear map
$\psi \circ \varphi : \mathbb{R}^d \to \mathbb{R}^m$ is represented in these bases by the matrix
$\Psi \cdot \Phi$. Show that
$$\|\Psi \cdot \Phi\| \leq \|\Psi\|\|\Phi\|,$$
where $\|\cdot\|$ is the Hilbert--Schmidt norm of a matrix (see \textit{Definition 9.96} in the
notes).

\textit{Official hint:} recall that $\|M\|^2$ is the sum of the squares of the entries of $M$.
Notice that $(\Psi\cdot\Phi)^i_j$ ($i$-th row, $j$-th column) is the scalar product of $\Psi^i$
and $\Phi_j$, which are vectors of $\mathbb{R}^n$. Apply Cauchy--Schwarz to each of them.
\end{exercise}

\begin{exercise}[3.11 --- (*) Uniform continuity and moduli of continuity]
\label{ex:3.11}
Let $(X,d_X)$ and $(Y,d_Y)$ be metric spaces and let $f : X \to Y$.

A \textbf{modulus of continuity} is a function $\omega : [0,\infty) \to [0,\infty]$ such that
$$\omega(0) = 0, \qquad \omega \text{ is nondecreasing}, \qquad \lim_{t \downarrow 0}\omega(t) = 0.$$
(Here $[0,\infty]$ means we also allow the value $+\infty$.) We say that $f$ \textbf{has modulus
of continuity} $\omega$ if
$$d_Y\big(f(x), f(y)\big) \leq \omega\big(d_X(x,y)\big) \qquad \forall x,y \in X.$$
\begin{enumerate}[label=\textbf{\arabic*.}]
  \item Prove that $f$ is uniformly continuous on $X$ if and only if $f$ has a modulus of
    continuity.
  \item \textit{(Bonus 1)} Assume $X \subset \mathbb{R}$ is an interval with the usual distance
    $d_X(x,y) = |x-y|$. Show that the modulus constructed in (1) can be chosen
    \textbf{subadditive}, i.e.\ $\omega(t_1+t_2) \leq \omega(t_1) + \omega(t_2)$ for all
    $t_1, t_2 \geq 0$.
  \item \textit{(Bonus 2)} Assume $X \subset \mathbb{R}^n$ is \textbf{convex}, meaning: for all
    $x_1, x_2 \in X$ and all $\lambda \in [0,1]$, $(1-\lambda)x_1 + \lambda x_2 \in X$. With
    $d_X(x,y) = \|x-y\|$ (Euclidean distance), show again that the modulus from (1) can be chosen
    subadditive.
\end{enumerate}

\begin{remark}[official]
In general metric spaces a modulus of continuity may have infinite value for finite $t > 0$.
Example: $X = \mathbb{Z}$ with the usual distance, $Y = \mathbb{R}$, $f(n) = n^2$. This $f$ is
uniformly continuous on $\mathbb{Z}$ (because points closer than 1 must be equal), but no finite
$\omega(1)$ can bound $|f(n+1) - f(n)| = 2n+1$ for all $n$. Allowing the value $+\infty$ avoids
this issue, and in many familiar settings (e.g.\ intervals/convex sets in $\mathbb{R}^n$) the
resulting $\omega_f(t)$ is finite for all $t$.
\end{remark}

\textit{Official hints:} \textbf{3.11.1} --- define the ``worst oscillation at scale $t$''
$\omega_f(t) := \sup\{d_Y(f(x),f(y)) : d_X(x,y) \leq t\}$. \textbf{3.11.2 / 3.11.3} --- if
$d_X(x,y) \leq t_1 + t_2$, choose an intermediate point $z$ on the segment from $x$ to $y$ so
that $d_X(x,z) \leq t_1$ and $d_X(z,y) \leq t_2$, then use the triangle inequality in $Y$.
\end{exercise}
